\section{Phát hiện bất thường trên dữ liệu log}

\subsection{Định nghĩa và phân loại}
Anomaly là sai lệch dữ liệu; fault là nguyên nhân tiềm ẩn; failure là sự kiện
vận hành có thể quan sát. Bất thường điểm liên quan tới một bản ghi, bất thường
tập thể liên quan tới chuỗi hoặc block, còn bất thường ngữ cảnh phụ thuộc
ngữ cảnh xung quanh.

\subsection{Các hướng tiếp cận}
Phương pháp giám sát cần nhãn train, bán giám sát chủ yếu học từ dữ liệu
bình thường, còn không giám sát không dùng nhãn huấn luyện. Với HDFS\_v1,
nhãn block cho phép giám sát ở mức block. BGL raw chưa gắn nhãn trong dữ liệu
hiện tại nên cần ground-truth riêng trước khi huấn luyện đánh giá giám sát.

\subsection{Học máy cổ điển}
KNN phát hiện trường hợp gần nhất theo không gian đặc trưng. Decision Tree
chia dữ liệu bằng các ngưỡng dễ kiểm tra. SLFN học biểu diễn phi tuyến nông.
LightAD là đại diện của họ này. Ưu điểm là chi phí thấp và triển khai đơn giản;
nhược điểm là phụ thuộc biểu diễn log và nhãn.

\subsection{Học sâu và Transformer}
LSTM mô hình hóa chuỗi event và quan hệ trước--sau. Transformer dùng attention
để học quan hệ dài hạn. LogFormer tiền huấn luyện Transformer rồi tinh chỉnh
adapter cho phát hiện bất thường. Hạn chế gồm chi phí GPU, dữ liệu huấn luyện
và độ nhạy parser.

\subsection{LLM kết hợp RAG}
LLM có thể đánh giá ngữ nghĩa log và giải thích quyết định. RAG truy xuất các
trường hợp liên quan để làm ngữ cảnh trước khi suy luận. LogRAG kết hợp
DeepSVDD cho ứng viên bất thường rồi dùng RAG/LLM tái đánh giá. Hạn chế gồm
độ trễ, chi phí API, tính không xác định và chất lượng retrieval.

\subsection{Chỉ số đánh giá}
Precision đo tỷ lệ cảnh báo đúng trong tổng cảnh báo. Recall đo tỷ lệ bất
thường được phát hiện. F1 cân bằng Precision và Recall. PR-AUC đặc biệt phù
hợp cho dữ liệu mất cân bằng. FPR đo tỷ lệ false alarm trên số mẫu bình
thường. Accuracy không nên là chỉ số chính cho HDFS\_v1/BGL.
